\section{Exact presentations on the three markings}
Apply the presentation-transport lemma of [U] to the three attached cores. Let $\beta=(d\ x)$ be the
bridge transposition and leave all tail transpositions unchanged.
For A and B take the core tree with transpositions
\[
 a=(12),\quad b=(23),\quad c=(34).
\]
Then $h_0=b$, $h_1=ac$, and the inverse words are
\begin{equation}\label{eq:abwords}
 a=(h_1\beta)^3,\qquad b=h_0,\qquad c=(h_1\beta)^3h_1.
                                                               \end{equation}
These identities hold for either $d=4$ or $d=3$: the bridge touches
one factor of the double transposition, and cubing its product with
$h_1$ isolates the other factor $(12)$.

For C take the core tree with
\[
 a=(12),\quad b=(13),\quad c=(34).
\]
The forward and inverse words are
\begin{equation}\label{eq:cwords}
 h_0=ac,\qquad h_1=babcba,\qquad
 a=(h_0\beta)^3,\quad b=(h_1\beta)^3,\quad c=(h_0\beta)^3h_0.
                                                               \end{equation}
Indeed $abcba=(24)$, so $babcba=(13)(24)$; the cube identities
isolate $(12)$ and $(13)$ respectively.
In each case the three core edges and the full tail form a tree on
$\Omega$. Substituting \eqref{eq:abwords} or \eqref{eq:cwords} into the tree-presentation theorem of [U] and
adding the two forward marking identities gives, by
the presentation-transport lemma of [U], a complete finite presentation on the
original $|\Omega|-2$ generators. This construction handles arbitrary
branching and specifies every relator by an explicit rule.

\subsection{The Coxeter diagrams and the printed local words}
Every pair order is determined by the following table, together with
incidence of edges in the tail:
\[
\begin{array}{c|ccc}
 &|h_0h_1|&|h_0\beta|&|h_1\beta|\\\hline
 A&4&2&6\\ B&4&3&6\\ C&2&6&6
\end{array}
\]
All other tail marks commute with both core marks. Two tail marks
have product order three if their tree edges meet, and order two
otherwise. In particular the bridge has order-three products with
\emph{every} edge at its other endpoint. To check the table, the two
D8 core products are four-cycles and the V4 core generators commute.
A transposition meeting just one factor of a double transposition
has a disjoint three-cycle and transposition as product, of order six.
The two remaining entries for $h_0\beta$ are the disjoint or
one-point-intersection cases for transpositions. This proves the
complete Coxeter-diagram description for every tail.

Let $t$ be any tail edge adjacent to the bridge at $x$. In A and B
write $a=(h_1\beta)^3=(12)$. Then $|ah_0|=3$ and $[a,t]=1$.
In A, $h_1\beta h_1=(3x)$, so $(h_0h_1\beta h_1)^3=1$.
In B, $h_1h_0h_1=(14)$ is disjoint from $\beta=(3x)$, so
$(h_1h_0h_1\beta)^2=1$. In C write
$a=(h_0\beta)^3=(12)$ and $b=(h_1\beta)^3=(13)$.
Then $|ab|=3$, $abh_1=(12)(24)$ has order three, and both
$a$ and $b$ commute with $t$. These are precisely the local
extra-word identities in [FL, Corollary 6.5(5)], in the displayed
coordinate conventions. Every triangle of tail edges also satisfies
the triangle relation of [U]. Thus the local diagram and word claims
are retained wherever their stated markings apply; the following
qualification concerns the uniform arbitrary-branching presentation.

\section{A literal branching qualification}
The source's accepted manuscript (institutional repository,
\texttt{10773/23861}, physical pages 17--18, Corollary 6.5) states
that $\rho_2$ commutes with all $\rho_i$ for $i\geq4$, while
Theorem 1.2 permits arbitrary branching immediately after the bridge.
These clauses cannot both hold literally. On nine letters, take
\begin{equation}\label{eq:branchingexample}
 (23),\ (12)(34),\ (45),\ (56),\ (57),\ (68),\ (79).
                                                               \end{equation}
The attachment theorem proves the full intersection property and
identifies the generated group as $S_9$. The complete flag proof also
applies. But $(45)(56)$ and $(45)(57)$ both have order three. Thus
the bridge fails to commute with two tail generators, independently
of how those tail labels are permuted. The double transposition and
its unique order-six neighbor identify the bridge within this marked
family, so a renumbering preserving the specified core cannot repair
that assertion.

The corrected general presentation is obtained from the transport
rule of [U] with \eqref{eq:abwords}--\eqref{eq:cwords}. For a path tail it specializes to the familiar changes of
marking used in the source. We do not assert, without a separate
Tietze calculation, that every redundant transported relation follows from
the source's printed shorter path list. Nor do we describe the false
branching clause as a recovered theorem. It is an explicit correction
within the intended general presentation result.

